% Central Paper C content (journal-neutral). Edit here; journal wrappers in ../journals/ only \input these files.
% Every number traces to ../../FINDINGS.md.
\section{Results}\label{sec:results}

\subsection{What open data contain}\label{sec:res-coverage}
In citing years up to 2014, 99.5\% of citation links have a citing item with a journal; this share falls to 69.3\%
in 2020 and recovers to 97.6\% in 2022 (Fig.~\ref{fig:coverage}). The dip is caused by one source. In 2020, 221,918 of
the 230,438 citing items typed \emph{dataset} without a journal (96.3\%) carry DOIs of the Global Biodiversity
Information Facility (prefix 10.15468, registered with DataCite): occurrence downloads whose reference lists cite the
datasets and taxonomic publications the records come from. These are legitimate data citations, but they are not the
journal citations Clarivate counts. Had they been counted, they would form 12.8\%, 14.4\% and 11.2\% of the JIF-window
numerator in 2019, 2020 and 2021, concentrated in a few taxonomy journals: in 2020, 99.8\% of the numerator of
\emph{Zootaxa} would come from such items, giving a JIF proxy above 600. Restricting citing items to those with a
journal removes the artifact; including them changes the JIF proxy noticeably only in 2019--2021 (correlation with the
main version 0.70--0.90, against at least 0.99 in other years).

\begin{figure}[t]
\centering
\includegraphics[width=0.6\linewidth]{F1_coverage_by_year.pdf}
\caption{Share of citation links whose citing item has a journal, whose both items have a journal, and whose both
items are journals under the journal rule, by citing year. The 2015--2021 dip is due to GBIF occurrence-download
records registered with DataCite.}
\label{fig:coverage}
\end{figure}

\subsection{Fidelity to Clarivate's published figures}\label{sec:res-fidelity}
Whether open data reproduce Clarivate's figures depends mainly on whether the citing journal's reference lists are in
OpenCitations (Fig.~\ref{fig:fidelity}). When they are, our self-citation share lies within $\pm10$ percentage points
of Clarivate's for 36 of 50 journals, with a median absolute difference of 4.6 points and a median signed difference
of $-2.8$ points. When they are absent, our share is zero and none of seven journals passes. For stacking pairs whose
donor has references and the pair has citations in our data (five pairs), the donor's share of the recipient's
numerator is reproduced for four and the exchange share for four. Many stacking donors in Clarivate's lists, however,
have no references in OpenCitations at all.

\begin{figure}[t]
\centering
\includegraphics[width=0.55\linewidth]{F2_fidelity.pdf}
\caption{Clarivate's published percentages against ours, by whether the citing journal's (self-citation) or donor's
(stacking) reference lists are present in OpenCitations. The band marks $\pm10$ percentage points.}
\label{fig:fidelity}
\end{figure}

The missing references are not random. For the journals whose references are absent, the median share of their DOIs
with deposited reference lists in Crossref is zero (zero for eight of eleven), against 0.96 for journals whose
references are present. The journals of one publisher that deposits references today began only recently (backfile
coverage around 5\%). These blind spots include donors in several Clarivate stacking cases.

On ranks, 59 of 85 self-citation suppressions (69\%) lie in the top 5\% of their year by self-citation share, with a
median percentile of 3.1\%, but only 9 are in the top 1\%. The criterion we fixed in advance for data year 2020,
at least half of the suppressed journals in the top 1\%, was not met (0 of 8); we report this as recorded.

\subsection{CIDRE on open data}\label{sec:res-cidre}
CIDRE finds between 11 and 32 groups per year, typically of three to five journals. It detects none of the 15 named
Clarivate pairs from 2018--2023 for which both journals are matched, in any of the eight partitions, and no recipient
or donor of these pairs is in any CIDRE group in its year. Some misses are due to data: the donor has no references in
OpenCitations, or the pair has too few citations to reach $\theta_w$. Two 2023 pairs, however, are well covered and
still missed.

For these, a stage-by-stage diagnosis shows the mechanism (Fig.~\ref{fig:peeling}). The stacking edge passes CIDRE's
test: 99.5\% (Optik to Ukrainian Journal of Physical Optics) and 85.5\% (Environmental Science and Pollution Research
to Climate Change Economics) of the donor's citations to the recipient fall in the JIF window, and the edges are
significant against the block model. The recipients' scores, 0.83 and 0.32, are well above $\theta$. But the donors are
large journals whose stacking citations are a small part of their output: their donor scores are 0.057 and 0.035.
Since CIDRE removes a journal when both of its scores are below $\theta$, the donor is removed, the recipient's score
falls to zero, and no group remains. The CIDRE paper motivates $\theta=0.15$ on the recipient side, noting that all
recipients suspended by Clarivate received at least 15\% of their citations from donors \citep{kojaku2021cidre}, but
applies the same threshold to donors.

\begin{figure}[t]
\centering
\includegraphics[width=0.6\linewidth]{F3_cidre_peeling.pdf}
\caption{Why CIDRE misses two well-covered 2023 stacking pairs: the recipients' scores exceed $\theta$, but the large
donors' scores do not, so the group-building step removes the donors.}
\label{fig:peeling}
\end{figure}

CIDRE does find stacking between small journals. In 2015 one of its groups consists exactly of Applied Clinical
Informatics and Methods of Information in Medicine, and in 2017 one consists exactly of Liver Cancer and Digestive
Diseases; in each case both journals were suppressed for stacking in that year, although Clarivate's spreadsheet for
those years does not name the pairs.

\subsection{The recipient-centric reading}\label{sec:res-variant}
Over 2015--2023, the variant flags 12 of 30 stacking-listed journals (40\%), against 3.4\% of all journals with at
least 20 numerator citations in some year: an enrichment of 11.7 ($p=1.2\times10^{-10}$). Without the two journals
used to design it, 10 of 28 are flagged (10.5-fold, $p=1.6\times10^{-8}$), and across partitions and size thresholds
the enrichment ranges from 5.7 to 10.5 ($p\le1.1\times10^{-5}$). For self-citation-listed journals the enrichment is
weaker and not robust across partitions (1.3 to 3.4), as expected for a measure of external citations.

Requiring flags to appear in more partitions lowers the base rate faster than the hit rate (Fig.~\ref{fig:enrichment}).
With flags present in all seven partitions, 2.3\% of journals are ever flagged, against 9 of 28 stacking-listed
journals: a 13.9-fold enrichment ($p=8.7\times10^{-9}$). Because the consensus rule was chosen after seeing these
results, we report all three levels.

\begin{figure}[t]
\centering
\includegraphics[width=0.6\linewidth]{F4_enrichment.pdf}
\caption{Share of listed journals ever flagged by the recipient-centric reading in 2015--2023, against the base rate,
by the number of partitions required (design cases excluded). Labels give the enrichment for stacking-listed journals.}
\label{fig:enrichment}
\end{figure}

Eight listed journals were flagged before their first listing (Fig.~\ref{fig:leadtime}). Granular Computing was
flagged in every year from 2016 to 2020, three to seven years before its 2023 listing, with Information Sciences, its
donor in Clarivate's list, as the largest excessive donor in four of those years. Panminerva Medica and Minerva Medica
were flagged three to four years before their 2023 listing, Liver Cancer one to two years before 2017 (largest
excessive donor Digestive Diseases), and Ukrainian Journal of Physical Optics in 2022 (donor Optik).

\begin{figure}[t]
\centering
\includegraphics[width=0.6\linewidth]{F5_lead_time.pdf}
\caption{Years in which listed journals were flagged (dots) and their first Clarivate listing (bars), partition
\emph{main}.}
\label{fig:leadtime}
\end{figure}

In the pre-specified evidence check on 196 randomly sampled consensus flags and 196 matched controls, 28.1\% of flagged
pairs carry at least one independent indicator against 6.1\% of controls ($p=6.2\times10^{-9}$, Fisher's exact
test). The most frequent indicator is a shared publisher of recipient and donor (19.4\% against 6.1\%); 5.1\% of
flagged pairs have a donor in Clarivate's lists and 5.6\% have strongly concentrated citations, against none of the
controls. These shares are a lower bound on the precision of the flags, not an estimate of it.

\subsection{Paper-level anatomy}\label{sec:res-anatomy}
Flagged pairs differ from matched controls in how their citations arise (Fig.~\ref{fig:anatomy}). Each citing paper
cites the recipient 1.98 times on median, against 1.28 ($p=2.5\times10^{-27}$, Mann--Whitney), and citations are more
unevenly spread over citing papers (Gini 0.25 against 0.18, $p=1.5\times10^{-8}$). The share from the top 10\% of
citing papers does not differ ($p=0.35$): the pattern is many papers each citing a little too often rather than a few
papers with long reference lists. Cited and citing titles are less similar in flagged pairs (TF--IDF cosine 0.084
against 0.133, $p=1.2\times10^{-24}$). Since excessive edges are defined against the block model, flagged pairs might
simply join more distant fields; in fact 27.9\% of flagged pairs share a block against 0.6\% of controls, and within
cross-block pairs the differences remain (cosine 0.084 against 0.133, $p=1.2\times10^{-17}$; citations per citing
paper 2.00 against 1.28). For the five named Clarivate pairs with citations the sample is too small for conclusions.

\begin{figure}[t]
\centering
\includegraphics[width=0.6\linewidth]{F7_anatomy.pdf}
\caption{Citations per citing paper (log scale) and citing--cited title similarity for named Clarivate pairs, flagged
pairs and their matched controls.}
\label{fig:anatomy}
\end{figure}

\subsection{Sensitivity to the block-model partition}\label{sec:res-stability}
With eight threads, block-model fits are not reproducible from a seed: six fits of the same network give 2,131 to 2,912
blocks, with pairwise normalized mutual information of 0.76--0.78 and adjusted Rand index \citep{hubert1985ari} of
0.44--0.54. Two single-threaded fits with the same seed give identical partitions, so the variation comes from the
parallel search. CIDRE's groups depend much less on the partition than the partitions themselves agree: the median
Jaccard similarity of group members between partitions is 0.85, and 93\% of groups have a counterpart (Fig.~\ref{fig:stability}).
The variant's flags are less stable (median Jaccard 0.62), which motivates the consensus rule.

\begin{figure}[t]
\centering
\includegraphics[width=0.6\linewidth]{F6_stability.pdf}
\caption{Agreement between fits: block-model partitions (ARI, NMI), CIDRE group members (Jaccard) and group recovery,
and the variant's flags (Jaccard).}
\label{fig:stability}
\end{figure}

